\documentclass[11pt,a4paper]{article}
\usepackage[margin=2.5cm]{geometry}

\usepackage[T1]{fontenc}   
\usepackage[utf8]{inputenc} 
\usepackage[english]{babel}

\usepackage{newtxtext,newtxmath}

\usepackage{microtype}

\usepackage{xcolor}
\definecolor{citationred}{RGB}{150,0,0}

\usepackage[
  colorlinks=true,
  citecolor=citationred,
  linkcolor=black,
  urlcolor=black
]{hyperref}

\usepackage[authoryear,round]{natbib}



\usepackage[most]{tcolorbox}  
\tcbuselibrary{skins}         


\newtcolorbox{commentbox}{
  colframe=green,        
  boxrule=0.4pt,         
  sharp corners,         
  enhanced,              
  left=2mm,right=2mm,    
  bottom=2mm,
  top=1.9em,             
  overlay={
    \node[anchor=north west,fill=citationred,text=white,
          font=\bfseries,inner ysep=2pt,inner xsep=4pt]
         at (frame.north west) {Commentaires};
  },
}

\renewcommand{\thefootnote}{\fnsymbol{footnote}}

\begin{document}


\section*{Introduction}

A modest reallocation of resources from richer to poorer countries would be enormous relative to the size of low-income economies. Transferring just 1\% of high-income countries' aggregate output to low-income countries (LICs) would be equivalent top double their national income\footnote{According to \href{https://data.worldbank.org/indicator/NY.GDP.PCAP.PP.CD?contextual=default&end=2024&locations=EU-ZG-XD-XM-1W-IN-US-CD-BI-LU-CN&start=2024&view=bar}{World Bank data for 2024}, GDP per capita in high-income countries (HICs) is 28 times higher than in low-income countries (LICs) at Purchasing Power Parity (PPP), and \href{https://data.worldbank.org/indicator/NY.GDP.PCAP.CD?end=2024&locations=EU-ZG-XD-XM-1W-IN-US-CD-BI-LU-CN&start=2024&view=bar}{68 times higher} in nominal terms. Approximately 625 million people live in the 25 low-income countries, defined here as countries with GDP per capita at PPP below \$1,135 per year, compared with 1.42 billion people living in high-income countries. On this basis, 1\% of HICs' GDP represents approximately 60\% of LICs' combined GDP using PPP values and 153\% using nominal values.}. A transfer of this magnitude would be more than sufficient to finance the eradication of extreme poverty \citep{sahoo_what_2025}.

The magnitude of these disparities, together with the potential impact of even limited redistribution, raises a broader question of global justice. In broad terms, global justice concerns the fair and sustainable distribution of socioeconomic and environmental resources both within and across countries, subject to the ecological constraints imposed by planetary boundaries. On this view, reducing extreme poverty and narrowing large international inequalities are necessary but insufficient objectives: distributive improvements cannot be considered sustainable if they are achieved through development paths that intensify climate change and expose already vulnerable populations to further environmental harm. Global justice therefore involves not only the distribution of income and wealth, but also access to public services, economic opportunities, environmental resources, and decision-making power.

International redistribution appears necessary not only to eradicate extreme poverty, a central objective of the universally endorsed Sustainable Development Goals, but also, over the longer run, to achieve convergence in GDP per capita across countries by 2100 \citep{fabre_JDS_2025}. The Global Justice Report estimates that this would require annual transfers from high-income to low-income countries equivalent to approximately 0.8\% of world GDP over the period 2026--2100 \citep[Figure~4.2]{bothe_et_al_2026}. The magnitude of existing global inequality underscores the potential role of such transfers. The richest 1\% of individuals in high-income countries receive more income than the entire bottom half of the global income distribution combined, while the richest 10\% receive more than half of global income \citep{chancel_world_2026}.

These reparative and institutional considerations are complemented by a distinct framework for international redistribution. Within an optimal-tax framework, \citet{kopczuk_limitations_2005} show that a centralized world income tax based on a border-neutral social welfare function would reduce the global Gini coefficient from 0.69 to 0.25, whereas decentralized national systems have almost no effect on global inequality. They conclude that efficiency costs alone cannot explain the limited scale of international redistribution, which is also constrained by free-riding among donor countries.

Moreover, the international distribution of resources cannot be understood solely through differences in current income or gross domestic product. Large flows of value continue to move from the Global South toward the Global North through unequal economic relations, debt servicing, tax-base erosion, and other forms of extraction \citep{hickel_plunder_2021}. Debt distress and austerity can constrain the policy choices available to low-income countries and reproduce asymmetric economic relationships \citep{fischer_return_2023}. The literature also emphasizes the historical origins of contemporary distributive obligations. Colonial institutions contributed to the formation of present economic structures and continue to influence development trajectories \citep{bhambra_colonial_2021}, while slavery and colonial extraction provide distinct grounds for reparative claims \citep{solbiac_britains_2015,darity_cumulative_2022}. Climate change adds a further dimension: because high-income countries have used a disproportionate share of the global carbon budget and, in particular, have already emitted beyond the levels compatible with the temperature limits set by the Paris Agreement, they may bear particular responsibilities toward countries that have contributed less to cumulative emissions but remain highly exposed to their consequences \citep{fanning_compensation_2023,Hickel_and_Zoomkawala_2026}. Taken together, contemporary inequality, historical responsibility, asymmetric economic relations, and climate liability provide several distinct arguments for transferring resources across national borders.

International redistribution should not, however, be conflated with all cross-border resource flows. Migrants' remittances, for example, are private household transfers rather than collectively determined redistributive policies. Officially recorded remittances to low- and middle-income countries reached approximately 0.61\% of world GDP in 2024, exceeding foreign direct investment and Official Development Assistance (ODA) combined \citep{ratha_remittances_2024}. Corporate profit shifting represents a different phenomenon: rather than a deliberate transfer toward poorer countries, it erodes national tax bases and limits governments' ability to finance public expenditure. Around 36\% of multinational profits are estimated to be shifted to tax havens, while associated global corporate-tax revenue losses have been estimated at 0.57\% and 0.74\% of world GDP each year,  although these estimates remain methodologically contested \citep{torslov_missing_2023, cobham_global_2018}. ODA is more directly redistributive and more fully institutionalized, but substantially smaller in aggregate terms. Development Assistance Committee donors provided 0.18\% of world GDP in ODA in 2024 on a grant-equivalent basis, followed by a preliminary decline to 0.15\% of world GDP in 2025 \citep{oecd_final_oda_2025, oecd_preliminary_oda_2026}. 
These comparisons indicate that the politics of international redistribution cannot be reduced to foreign aid alone. It also concerns international taxation, fiscal sovereignty, debt, climate finance, and the institutional rules governing cross-border resource allocation.

This chapter does not seek to resolve the normative question of how much international redistribution justice ultimately requires. It addresses a narrower positive question: to what extent citizens are prepared to accept the policies and institutions through which international redistribution could be implemented? The available literature primarily examines respondents in high-income or net-contributing countries, although several studies provide broader cross-national or globally representative evidence. The chapter therefore focuses principally on the attitudes of citizens  believe that they will bear, part of the costs of international redistribution.

The chapter starts from precisely this puzzle. International redistribution appears normatively important and economically relevant, yet it remains rarely implemented. The argument developed here is a pattern of conditional acceptance. Support tends to be greater when policies are perceived as fair, effective, fiscally credible, progressively financed, and embedded in plausible arrangements for international burden-sharing. Conversely, approval may remain politically weak when citizens doubt that resources will reach intended beneficiaries, perceive foreign expenditure as competing with domestic redistribution, or are held back by the costs involved.

In this chapter, we develop this argument in seven steps. We first examine moral universalism as a psychological and political foundation for extending distributive concern beyond national borders. We then turn to foreign aid, the most familiar public instrument of international redistribution, and studies 
the level of support and how support depends on information, perceived effectiveness, institutional credibility, ideology, and competition with domestic expenditure. In the third section, we consider rule-based forms of redistribution, particularly international taxation of rich individuals funding poverty reduction and public services in low-income countries.
We analyse then moves from instruments to institutions. In the fourth section, we examine whether citizens accept the forms of global governance required to coordinate contributions, constrain free-riding, and implement policies across countries. 

The fifth section turns to international carbon pricing, which would make ordinary citizens in high-income countries bear the costs, and not only the richest people. The sixth section distinguishes collective policy approval from stated readiness to contribute personally, paying particular attention to the gap between respondents’ own propensity to contribute and their beliefs about the propensity of others to do so. 
Finally, the chapter considers why internationally redistributive and climate policies remain politically marginal despite their conditional acceptability. It examines low issue salience, factual misperceptions, second-order beliefs, reference groups, fiscal competition, biased perceptions among political representatives, and weaknesses in political supply.

\newpage
\bibliographystyle{plainnat}
\bibliography{chapter/references}

\end{document}